\documentclass[11pt,a4paper]{article}
\usepackage[utf8]{inputenc}
\usepackage[margin=1in]{geometry}
\usepackage{hyperref}
\usepackage{graphicx}
\usepackage{amsmath}
\usepackage{booktabs}

\title{RSCT Review: Boosting deep Reinforcement Learning using pretraining with Logical Options}
\author{Swarm-It Research Discovery\\\small Automated RSCT-Based Analysis}
\date{March 09, 2026}

\begin{document}
\maketitle

\begin{abstract}
This document provides an automated review of the paper ``Boosting deep Reinforcement Learning using pretraining with Logical Options'' from an RSCT (Representation-Solver Compatibility Testing) perspective. The paper achieved an RSCT relevance score of 36.7% and topic similarity of 43.0%.
\end{abstract}

\section{Paper Information}
\begin{itemize}
\item \textbf{Title:} Boosting deep Reinforcement Learning using pretraining with Logical Options
\item \textbf{Authors:} Zihan Ye, Phil Chau, Raban Emunds, Jannis Blüml, Cedric Derstroff et al.
\item \textbf{Source:} \url{https://arxiv.org/abs/2603.06565v1}
\item \textbf{RSCT Relevance:} 36.7%
\item \textbf{Topic Match:} 43.0%
\item \textbf{Key Concepts:} N/A
\end{itemize}

\section{Original Abstract}
Deep reinforcement learning agents are often misaligned, as they over-exploit early reward signals. Recently, several symbolic approaches have addressed these challenges by encoding sparse objectives along with aligned plans. However, purely symbolic architectures are complex to scale and difficult to apply to continuous settings. Hence, we propose a hybrid approach, inspired by humans' ability to acquire new skills. We use a two-stage framework that injects symbolic structure into neural-based reinforcement learning agents without sacrificing the expressivity of deep policies. Our method, called Hybrid Hierarchical RL (H\textasciicircum{}2RL), introduces a logical option-based pretraining strategy to steer the learning policy away from short-term reward loops and toward goal-directed behavior while allowing the final policy to be refined via standard environment interaction. Empirically, we show that this approach consistently improves long-horizon decision-making and yields agents that outperform strong neural, symbolic, and neuro-symbolic baselines.

\section{RSCT Analysis}
This paper presents research with 37% relevance to RSCT theory.

\textbf{Abstract Summary:} Deep reinforcement learning agents are often misaligned, as they over-exploit early reward signals. Recently, several symbolic approaches have addressed these challenges by encoding sparse objectives along with aligned plans. However, purely symbolic architectures are complex to scale and difficult to apply to continuous settings. Hence, we propose a hybrid approach, inspired by humans' ability to acquire new skills. We use a two-stage framework that injects symbolic structure into neural-based ...

\textbf{RSCT Relevance:} The work touches on concepts related to general machine learning, which have connections to the Representation-Solver Compatibility Testing framework.

\textbf{Further Analysis:} A detailed manual review is recommended to fully assess the implications for RSCT-based certification approaches.


\subsection*{RSCT Certification Metrics}
\begin{tabular}{ll}
\textbf{Relevance (R):} & 0.374 \\
\textbf{Spurious (S):} & 0.319 \\
\textbf{Noise (N):} & 0.307 \\
\textbf{Kappa ($\kappa$):} & 0.549 \\
\end{tabular}

The kappa score of 0.549 indicates moderate representation-solver compatibility.


\section{Relevance to Swarm-It}
This paper was identified by the Swarm-It research discovery pipeline as potentially relevant to RSCT-based AI certification. The combined score of 39.2% places it in the moderate relevance category for further investigation.

\vspace{1em}
\hrule
\vspace{0.5em}
\small{Generated by Swarm-It Research Discovery | \url{https://swarms.network} | RSCT Certified}

\end{document}
